\section{Introduction}
\label{sec:introduction}

Transformer language models compute by applying a stack of residual blocks, and
every token pays for every block. Depth is therefore a direct axis of inference
cost: it multiplies latency, memory traffic, and energy, and unlike hidden width
it cannot be amortized away by batching. As models grow deeper to grow more
capable, the cost of simply traversing the stack has become one of the dominant
terms in the inference bill \citep{zhen2025taming}.

Layer skipping is an attractive response. Removing an entire block eliminates
real computation --- its attention and MLP matmuls both disappear --- without
changing the hidden width, the attention pattern, or the KV-cache layout, so it
composes with the rest of the inference stack for free. A substantial literature
establishes that Transformers tolerate this surprisingly well: blocks in large
models are highly redundant, and a well-chosen subset can be deleted with modest
degradation \citep{men2024shortgpt, gromov2024unreasonable, fan2020reducing}.
Almost all of this work asks the same question --- \emph{which} blocks to remove
--- and answers it with progressively better importance scores. Far less attention
has gone to the operator that does the removing. When a block is skipped, its
residual update is replaced by zero. That choice is never argued for; it is simply
what ``removing a layer'' has come to mean. It is computationally convenient, and
representationally wasteful.

A separate line of work, meanwhile, tells us that the residual stream is not a
sequence of unrelated states but a trajectory. Residual networks behave like
ensembles of shallow paths whose contributions accumulate smoothly
\citep{veit2016residual}, residual connections encourage successive blocks to
iteratively refine a single evolving representation
\citep{jastrzebski2018residual}, and a hidden state at one depth can often be
mapped to a nearby depth by a linear transformation \citep{din2023jump,
belrose2023eliciting}. If depth is a trajectory, a block's update is not an
isolated quantity: it is the next step along a path the model has already been
walking, and steps along a smooth path should be predictable from the steps
before them.

\begin{center}
\emph{How much of a skipped Transformer block's residual update \\ is predictable
from the preceding depth trajectory, \\ and can it be recovered at negligible
cost?}
\end{center}

We take this question literally. Writing the pre-norm update at layer $\ell$ as
$\dd_\ell = F_\ell(\hh_\ell)$, so that $\hh_{\ell+1} = \hh_\ell + \dd_\ell$,
plain skipping is precisely the estimator $\widehat{\dd}_\ell = 0$ --- a
zero-order predictor that ignores every update the model has already computed. We
instead treat $\dd_{\ell-1}, \dd_{\ell-2}, \dots$ as a short series indexed by
depth and fit an autoregressive model that forecasts $\dd_\ell$ from it, using
information already sitting in the residual stream at inference time, free of
charge.

The answer to the question is more interesting than a clean yes. Depth is
\emph{not} the smooth forward trajectory the analogy suggests. Through the middle
of Qwen2.5-0.5B, consecutive residual updates are mildly \emph{anti}-correlated:
every fitted coefficient is negative across layers 4--16, and \phm{88}\% of the
individual channels agree on that sign, so successive blocks there partially undo
one another rather than continuing a path. Genuine depth-wise momentum exists, but
it is confined to the final quarter of the stack. And the property that looks like
the natural way to decide \emph{where} to skip --- how well a layer's update can be
forecast --- turns out to be almost useless for that purpose: predictability and
recoverability are statistically unrelated (\Cref{sec:experiments}).

We introduce \textbf{\mname (Depth-wise AutoRegressive update prediction)}, which
replaces the zero-update assumption with a closed-form per-channel forecast of the
missing update, fitted on unlabeled text with no gradient training. \mname
contributes three results. (1) \textbf{The zero-update assumption is unnecessarily
destructive, but only a per-channel predictor exploits that.} A single global
coefficient per layer --- the obvious first instantiation --- fails almost
everywhere, and at some layers is actively harmful: at layer 3 of Qwen2.5-0.5B,
substituting the scalar's prediction is markedly \emph{worse} than dropping the
update outright. Granting the predictor one coefficient per channel, and relaxing
nothing else, is what separates a method that works from one that does not; it
costs $d$ scalars per skipped block and leaves inference at $O(Td)$. (2)
\textbf{Predictability is not recoverability.} Held-out explained update energy
$P_\ell$ --- the cheap, natural selection signal, computable from the dense model
before any layer is skipped --- is uncorrelated with the damage a predictor
actually repairs (Spearman \phm{$-0.015$}, $p=\phm{0.95}$, $n=\phm{22}$). Selecting
the most predictable layers is not merely suboptimal but actively bad: at two
skipped blocks it leaves the scalar predictor far \emph{worse} than plain skipping.
Layer selection must be validated against recovery, which must be measured. (3)
\textbf{Likelihood recovery does not transfer to behaviour.} \mname recovers a
real fraction of the language-modelling damage that skipping causes at every scale
and budget we test --- \phm{10.3}\% at the deployable operating point of
Qwen2.5-1.5B, and more when layers are selected to maximise it. Downstream
accuracy does not follow. At four skipped blocks, at the same operating point, the
accuracy gap is not closed but slightly widened: \phm{$-21.1$}\% on Qwen2.5-0.5B
and \phm{$-6.0$}\% on Qwen2.5-1.5B, against positive likelihood recovery in both
cases. We therefore make no claim that \mname improves downstream quality, and we
report the one budget at which it does (\phm{18.5}\% at $k=2$ on 1.5B) as the
exception it is.

The three findings are one finding. Predictability is an energy-weighted geometric
score, and it does not predict repair. The ridge coefficient chosen by that same
geometric criterion leaves the model catastrophically worse. And the likelihood
recovery \mname does achieve does not reach the behaviour anyone actually deploys
a model for. Each is a cheaper proxy standing in for functional preservation, and
each overestimates it.
